\documentclass[10pt,letterpaper]{article} % exam class defaults to 10pt on letterpaper
\usepackage[margin=1in]{geometry}         % exam class defaults to 1-inch margins
\usepackage{graphicx} 
\usepackage{fancyhdr}
\usepackage{amssymb}
\usepackage{amsmath}
\usepackage{mathtools}
\usepackage[most]{tcolorbox}
\usepackage[dvipsnames]{xcolor}
\usepackage{blindtext}
\usepackage{amsthm}
\usepackage{upgreek}
\usepackage{hyperref}
\usepackage{tikz}
\usetikzlibrary{positioning}
%making an edit, want to see if auto script commits it%
%making another edit to test the commit bot

\setlength{\parindent}{0pt}
\setlength{\parskip}{\baselineskip} % Matches the paragraph spacing of exam
\hbadness=10000 % Place in your preamble


\newtheorem{thm}{Theorem}[section]
\newtheorem{lem}{Lemma}[section]
\newtheorem{defn}{Definition}[section]
\newtheorem{cor}{Corollary}[section]
\newtheorem{axm}{Axiom}[section]

\tcolorboxenvironment{thm}{colback=lime, colframe=black, breakable}
\tcolorboxenvironment{lem}{colback=yellow, colframe=black, breakable}
\tcolorboxenvironment{defn}{colback=pink, colframe=black, breakable}
\tcolorboxenvironment{cor}{colback=SpringGreen, colframe=black, breakable}
\tcolorboxenvironment{axm}{colback=white, colframe=black, breakable}

% --- Header/Footer Setup ---
\fancyhf{}
\renewcommand\headrulewidth{0.5pt}
\pagestyle{fancy}
\author{Omkar Tasgaonkar}
\setlength{\headsep}{0.5in}
\fancyfoot[C]{(\thepage)}
\renewcommand\footrulewidth{0.5pt}

\title{Probability Theory}
\author{Omkar Tasgaonkar \\ \small otasgaonkar@ucla.edu}
\date{}

\fancypagestyle{firstpage}{
    \fancyhf{}
    \fancyhead[C]{\textbf{Notes in Analysis}}
    \renewcommand\headrulewidth{0.5pt}
    \fancyfoot[C]{(\thepage)}
    \renewcommand\footrulewidth{0.5pt}
}
\begin{document}
\maketitle
These are my notes on probability theory, which I am studying for C271A and as an extension of what I learned in analysis.
I hope to cover the relevant measure theory whenever possible, and I will also add completed practice problems as well.

\begin{tcolorbox}[colback=white, colframe=black, arc=0mm]
Textbooks Used to Develop these Notes:
\begin{enumerate}
\item \textit{\textbf{Probability: Theory and Examples} - Rick Durrett}
\item \textit{\textbf{Real Analysis} - Gerald Folland}
\end{enumerate}
Other Textbooks Referenced:
% \begin{enumerate}
% \item \textit{\textbf{Textbook} - Author}
% \end{enumerate}
\end{tcolorbox}
\begin{tcolorbox}[colback=white, colframe=black, arc=0mm]
\textbf{Notes Last Updated:} \today
\end{tcolorbox}
\thispagestyle{firstpage}
\newpage
\begin{tcolorbox}[colback=white, colframe=black, arc=0mm]
\tableofcontents
\end{tcolorbox}
\textbf{NOTE:} Any sections/subsections marked with an asterisk are extra sections and have been added to this document as extensions of introduced theory.
\pagebreak

\section{Probability Spaces}
\subsection{Algebras and Rings}
We begin with some basic measure theory which will motivate the formulation of the $\sigma$-algebra and allow us to go from generality to specificity in defining probability spaces.
%I will get back to this at a later time

We begin by defining a probability space as a triple $(\Omega, \mathcal{F}, P)$. The first element $\Omega$ is called the set of outcomes or the "sample space", which is just an underlying set that can be defined as we wish. 
\begin{defn}[$\sigma$-algebra]
The second element $\mathcal{F}$ is called a $\sigma$-algebra, satisfying the following properties:
\begin{enumerate}
\item $\Omega \in \mathcal{F}$
\item $\forall S, P \in \mathcal{F}: S \setminus P \in \mathcal{F}$
\item The countable (finite or countably infinite) union of sets in $\mathcal{F}$ is in $\mathcal{F}$:
\begin{equation*}
\forall \{S_n\}_{n \in \mathbb{N}} \in \mathcal{F}^\mathbb{N}: \bigcup_{n \in \mathbb{N}}S_n \in \mathcal{F}
\end{equation*}
We can equivalently replace this condition with closure under disjoint unions.
\end{enumerate}
We call the set $S$ an event (specifically an event with an assigned probability, which we shall see later).

We also define an algebra over sets $\mathcal{A}$, which weakens condition 3 to finite unions.
\end{defn}

\begin{lem}
We can derive other properties of sets in the $\sigma$-algebra, such as:
\begin{enumerate}
\item $\varnothing \in \mathcal{F}$
\item $\forall \{S_n\}_{n \in \mathbb{N}} \in \mathcal{F}^\mathbb{N}$:
\begin{equation*}
\bigcap_{n \in \mathbb{N}} S_n \in \mathcal{F}
\end{equation*}
\item $\forall S \in \mathcal{F}: \Omega \setminus S \in \mathcal{F}$
\end{enumerate}
These properties also hold for algebras in the finite sense.
\end{lem}
\begin{proof}
\textbf{$\sigma$-algebras}: For $\varnothing \in \mathcal{F}$, this follows from the fact that $\Omega \in \mathcal{F}$ and from property two, $\Omega \setminus \Omega = \varnothing \in \mathcal{F}$.

The second statement follows from DeMorgan's Laws. If $S_n \in \mathcal{F}$, then $S_n^c \in \mathcal{F}$, which tells us that:
\begin{equation*}
\bigcup_{n \in \mathbb{N}} S_n^c = \left(\bigcap_{n \in \mathbb{N}}S_n\right)^c \in \mathcal{F}
\end{equation*}
Once again by the completement axiom, we see that:
\begin{equation*}
\bigcap_{n \in \mathbb{N}}S_n \in \mathcal{F}
\end{equation*}

As for the third statement, we know that since $\Omega \in \mathcal{F}$, then we have that for all $S \in \mathcal{F}$, that $\Omega \setminus S \in \mathcal{F}$. 
\end{proof}

\begin{proof}
\textbf{Algebras}: $\varnothing \in \mathcal{F}$ follows from property 2.

Since $\Omega \in \mathcal{F}$, we have that $\forall S \in \mathcal{F}: \Omega \setminus S \in \mathcal{F}$. So, for any finite collection $\{S_i\}_{i = 0}^n$, we have that $S_i^c \in \mathcal{F}$:
\begin{equation*}
\bigcup_{i=0}^n S_i^c \in \mathcal{F} \implies \left(\bigcup_{i=0}^n S_i^c\right)^c = \bigcap_{i=0}^n S_i \in \mathcal{F}
\end{equation*}
\end{proof}

\begin{lem}
The intersection of a family of $\sigma$-algebras for a set $\Omega$ is a $\sigma$-algebra.
\end{lem}
\begin{proof}
Let us call the family $\{\mathcal{F}_\alpha: \alpha \in I\}$ such that we have valid measurable spaces $(\Omega, \mathcal{F}_\alpha)$, then:
\begin{equation*}
\mathcal{F} \coloneq \bigcap_{\alpha \in I}\mathcal{F}_\alpha
\end{equation*}
is a valid sigma algebra such that $(\Omega, \mathcal{F})$ is a measurable space. We see that:
\begin{equation*}
\forall \alpha \in I: \Omega \in \mathcal{F}_\alpha \implies \Omega \in \bigcap_{\alpha \in I}\mathcal{F}_\alpha
\end{equation*}
so property $1$ in Definition 1.1 is fulfilled. Now, for some $S$ in the intersection:
\begin{equation*}
S \in \bigcap_{\alpha \in I}\mathcal{F}_\alpha \iff \forall \alpha \in I: S \in \mathcal{F}_\alpha
\end{equation*}
\begin{equation*}
\implies \forall \alpha \in I: S^c \in \mathcal{F}_\alpha \iff S^c \in \bigcap_{\alpha \in I}\mathcal{F}_\alpha
\end{equation*}
thus property $2$ is fulfilled. Now consider some $\{S_n\}_{n \in \mathbb{N}} \in \mathcal{F}^\mathbb{N}$, then:
\begin{equation*}
\forall n \in \mathbb{N}: S_n \in \mathcal{F} \implies \forall \alpha \in I: S_n \in \mathcal{F}_\alpha 
\end{equation*}
\begin{equation*}
\implies \forall \alpha \in I: \{S_n\}_{n \in \mathbb{N}} \in \mathcal{F}_\alpha^\mathbb{N} \implies \forall \alpha \in I: \bigcup_{n \in \mathbb{N}}S_n \in \mathcal{F}_\alpha
\end{equation*}
\begin{equation*}
\implies \bigcup_{n \in \mathbb{N}}S_n \in \mathcal{F}
\end{equation*}
so we have verified property $3$. This is quite similar to topologies, where the intersection of topologies is a topology, but the union is not necessarily a topology.
\end{proof}

Consider some sample space $\Omega \coloneq \{1, 2, 3, 4\}$, then consider:
\begin{equation*}
\mathcal{F}_1 \coloneq \{\Omega, \varnothing, \{1\}, \{2, 3, 4\}\}
\end{equation*}
\begin{equation*}
\mathcal{F}_2 \coloneq \{\Omega, \varnothing, \{2\}, \{1, 3, 4\}\}
\end{equation*}
both of which are $\sigma$-algebras. If we take the intersection:
\begin{equation*}
\mathcal{F}_3 \coloneq \mathcal{F}_1 \cap \mathcal{F}_2 = \{\Omega, \varnothing\}
\end{equation*}
which is the "coarsest" or trivial $\sigma$-algebra. Now taking the union:
\begin{equation*}
\mathcal{F}_1 \cup \mathcal{F}_2 = \{\Omega, \varnothing, \{1\}, \{2\}, \{2, 3, 4\}, \{1, 3, 4\}\}
\end{equation*}
but by Definition 1.1 property $3$, we would need $\{1, 2\}$ to belong in the union, thus the union is not necessarily a $\sigma$-algebra.

An example of a $\sigma$-algebra is the countable $\sigma$-algebra, defined as:
\begin{equation*}
\mathcal{F} \coloneq \{S \subseteq \Omega: S \ \text{countable} \lor S^c \ \text{countable}\}
\end{equation*}
\begin{proof}
We have that $\Omega \in \mathcal{F}$ since $\varnothing$ is trivially countable. Now consider some $S \in \mathcal{F}$, then if $S$ is countable, we know that
$(S^c)^c$ is countable, thus $S^c \in \mathcal{F}$. If $S$ is not countable, then we know that $S^c$ must be countable, so $S^c \in \mathcal{F}$. 

Now take a countable union of sets in the sequence $\{S_n\}_{n \in \mathbb{N}} \in \mathcal{F}^\mathbb{N}$. If each $S_n$ is countable, then it follows that the union is in $\mathcal{F}$ from the fact
that countable union of countable sets is countable. If each $S_n$ is not countable, then:
\begin{equation*}
\left(\bigcup_{n \in \mathbb{N}}S_n\right)^c = \bigcap_{n \in \mathbb{N}}S_n^c \in \mathcal{F}
\end{equation*}
since $S_n$ not countable and $S_n \in \mathcal{F}$ means $S_n^c$ is countable. Assume some are countable and some are not, then for the $S_i$ that are not countable, we have that $S_i^c$ is countable. This forces the complement of the union to be countable.
\end{proof}

\begin{defn}
Given a space $\Omega$ and some set $S \subseteq \mathcal{P}(\Omega)$, we define the minimal algebra generated by $S$ as:
\begin{equation*}
\mathcal{F}(S) \coloneq \bigcap_{S \subseteq \mathcal{F}}\mathcal{F}
\end{equation*}
which is an algebra as per Lemma 1.2.
\end{defn}
Consider two subsets $P \subseteq S \subseteq \mathcal{P}(\Omega)$, then we know that:
\begin{equation*}
P \subseteq S \subseteq \mathcal{F}(S) \implies \mathcal{F}(P) \subseteq \mathcal{F}(S)
\end{equation*}
Since $\mathcal{F}(S)$ is a $\sigma$-algebra containing $P$ thus belongs in the intersection of the minimal $\sigma$-algebra containing $P$. We say the Borel $\sigma$-algebra on $\mathbb{R}$ is one that is generated by the set of all open intervals:
\begin{equation*}
I \coloneq \{(a, b): (\forall a, b \in \mathbb{R}: a < b)\} \subseteq \mathcal{P}(\mathbb{R})
\end{equation*}
\begin{equation*}
\mathcal{B} \coloneq \mathcal{F}(I)
\end{equation*}
The motivation of defining $\mathcal{B}$ as a minimal algebra of an algebra generated by open intervals is to essentially "hope" that unmeasurable/pathological sets like the Vitali set are eliminated from question. $\mathcal{B}$ will contain open intervals, closed intervals, and countable unions and intersections of them (so Cantor type sets are also included).

The issue is, while the Borel $\sigma$-algebra gives us candidates for what sets are measurable, it does not give the existence of a measure in $\mathbb{R}$.

Consider the minimal $\sigma$-algebra of an already existing $\sigma$-algebra $\mathcal{G}$, with the minimal $\sigma$-algebra being $\mathcal{F}(\mathcal{G})$. Then:
\begin{equation*}
\mathcal{F}(\mathcal{G}) \coloneq \bigcap_{\mathcal{G} \subseteq \mathcal{F}}\mathcal{F}
\end{equation*}
By definition, we have that $\mathcal{G} \subseteq \mathcal{F}(\mathcal{G})$. However, since $\mathcal{G} \subseteq \mathcal{G}$, then we have that:
\begin{equation*}
\mathcal{F}(\mathcal{G}) \subseteq \mathcal{G} \implies \mathcal{F}(\mathcal{G}) = \mathcal{G}
\end{equation*}

We introduce a few more concepts:
\begin{defn}
A ring $\mathcal{R} \subseteq \mathcal{P}(\Omega)$ is an algebera without the condition that $\Omega \in \mathcal{R}$. A $\sigma$-ring $\mathcal{R}_\sigma$ is a $\sigma$-algebra without the same condition.
\end{defn}

\begin{lem}
Properties of rings:
\begin{enumerate}
\item Rings (and respectively $\sigma$-rings) are closed under (countable) intersections
\end{enumerate}
\end{lem}
\begin{proof}
\textbf{Property One:} Consider a ring $\mathcal{R}$ and a finite collection of sets $\{S_i\}_{i = 0}^n \in \mathcal{R}^{n+1}$.
We prove by induction, consider the base case $n=1$. Then:
\begin{equation*}
S_1 \cap S_0 = S_1\setminus(S_1 \setminus S_0)
\end{equation*}
Since $S_1, S_0 \in \mathcal{R}$, then $S_1 \setminus S_0 \in \mathcal{R}$ and applying the set difference property again, we obtain that the intersection is in the ring. Now
assume this to be true for some $n \in \mathbb{N}$. Then for $n+1$:
\begin{equation*}
\bigcap_{i=0}^{n+1}S_i = S_{n+1} \cap \bigcap_{i=0}^n S_i = S_{n+1} \setminus \left(S_{n+1} \setminus \bigcap_{i=0}^n S_i\right) \in \mathcal{R}
\end{equation*}
Now consider a $\sigma$-ring $\mathcal{R}_\sigma$ and a sequence of sets $\{S_n\}_{n \in \mathbb{N}} \in \mathcal{R}_\sigma^\mathbb{N}$. We know that:
\begin{equation*}
\bigcup_{n \in \mathbb{N}}S_n \in \mathcal{R}_\sigma
\end{equation*}
Now by the finite set difference property, we know that $E\setminus S_n \in \mathcal{R}_\sigma$, so their countable union is in the ring too. Now:
\begin{equation*}
E \setminus \bigcup_{n \in \mathbb{N}}(E \setminus S_n) = E \cap \left(\bigcup_{n \in \mathbb{N}}(E \setminus S_n)\right)^c = E \cap \bigcap_{n \in \mathbb{N}}(E \setminus S_n)^c = E \cap \bigcap_{n \in \mathbb{N}}S_n
\end{equation*}
If we set $E = S_0$, then:
\begin{equation*}
S_0 \cap \bigcap_{n \in \mathbb{N}}S_n = \bigcap_{n \in \mathbb{N}}S_n \in \mathbb{R}_\sigma
\end{equation*}
by once again applying the set difference property to $E$ setminus the countable union (in the case that $E = S_0$).
\end{proof}

\subsubsection{Problems (Folland and Durrett)}
\textbf{Folland 1.2.1D}: If $\mathcal{R}_\sigma$ is a $\sigma$-ring, then $\mathcal{F} \coloneq \{E \subseteq X : (\forall F \in \mathcal{R}_\sigma: E \cap F \in \mathcal{R}_\sigma)\}$ is a $\sigma$-algebra:
\begin{proof}
First we observe that $X$ is in $\mathcal{F}$ since for all $F \in \mathcal{R}_\sigma$, we have that $X \cap F = F \in \mathcal{R}_\sigma$. Now, consider any $E, G \in \mathcal{F}$ such that:
\begin{equation*}
\forall F \in \mathcal{R}_\sigma: E \cap F \in \mathcal{R}_\sigma \land G \cap F \in \mathcal{R}_\sigma
\end{equation*}
Then by set difference properties of the $\sigma$-ring, we have that:
\begin{equation*}
(E \cap F) \setminus (G \cap F) \in \mathcal{R}_\sigma \implies (E \cap F) \cap (G^c \cup F^c) = E \cap G^c \cap F = (E \setminus G) \cap F \in \mathcal{R}_\sigma
\end{equation*}
thus $E\setminus G \in \mathcal{F}$. Finally, we must prove closure under countable union. Take a sequence $\{E_n\}_{n \in \mathbb{N}} \in \mathcal{F}^\mathbb{N}$ such that:
\begin{equation*}
\forall F \in \mathcal{R}_\sigma, \forall n \in \mathbb{N}: E_n \cap F \in \mathcal{R}_\sigma
\end{equation*}
\begin{equation*}
F \cap \bigcup_{n \in \mathbb{N}}E_n = \bigcup_{n \in \mathbb{N}}(F \cap E_n)
\end{equation*}
Since $F \cap E_n \in \mathcal{R}_\sigma$ for each $E_n$, then by closure under countable union of the $\sigma$-ring, we have that the countable union of $E_n$ is in $\mathcal{F}$.
Thus $\mathcal{F}$ is a $\sigma$-algebra. Here we used a $\sigma$-ring to describe a $\sigma$-algebra, a similar technique to what is used in Caratheodory's Extension Theorem.
\end{proof}

\textbf{Folland 1.2.5}: If $\mathcal{M}$ is the $\sigma$-algebra generated by $E \subseteq \mathcal{P}(X)$, then $\mathcal{M}$ is the union of the $\sigma$-algebras generated by $\mathcal{S}$ as $\mathcal{S}$ ranges over all countable subsets of $E$. In other words, we want to show:
\begin{equation*}
\mathcal{F}(E) = \bigcup_{\alpha \in I}\mathcal{F}(S_\alpha)
\end{equation*}

\begin{proof}
Since $E \subseteq \mathcal{F}(E)$, we have that:
\begin{equation*}
\forall \alpha \in I: S_\alpha \subseteq E \implies \forall \alpha \in I: \mathcal{F}(S_\alpha) \subseteq \mathcal{F}(E) \implies \bigcup_{\alpha \in I}\mathcal{F}(S_\alpha) \subseteq \mathcal{F}(E)
\end{equation*}
For the reverse inclusion, we show that the union is a sigma algebra itself. The space $X$ is trivially in the union. If we take some $P$ in the union, then:
\begin{equation*}
\exists \alpha \in I: P \in \mathcal{F}(S_\alpha) \implies P^c \in \mathcal{F}(S_\alpha) \implies P^c \in \bigcup_{\alpha \in I}\mathcal{F}(S_\alpha)
\end{equation*}
Now consider some sequence $\{P_n\}_{n \in \mathbb{N}}$ made from elements in the union. Then we have that:
\begin{equation*}
\forall n \in \mathbb{N}, \exists \alpha \in I: P_n \in \mathcal{F}(S_\alpha)
\end{equation*}
Denote each of these $\alpha$ as $\alpha_n$ based on the corresponding set $P_n$. Then we have that:
\begin{equation*}
\bigcup_{n \in \mathbb{N}}S_{\alpha_n}
\end{equation*}
is a countable union of countable subsets of $E$ and thus is a countable subset of $E$, so there must exist some $\alpha'$ such that:
\begin{equation*}
\forall n \in \mathbb{N}: S_{\alpha_n} \subseteq \bigcup_{n \in \mathbb{N}}S_{\alpha_n} = S_{\alpha'} \implies \mathcal{F}(S_{\alpha_n}) \subseteq \mathcal{F}(S_{\alpha'})
\end{equation*}
Thus there exists an $\alpha' \in I$ such that:
\begin{equation*}
\forall n \in \mathbb{N}: P_n \in \mathcal{F}(S_{\alpha'}) \subseteq \bigcup_{\alpha \in I}\mathcal{F}(S_\alpha)
\end{equation*}
so the union is a valid $\sigma$-algebra. Now if we define $E \coloneq \{E_\beta : \beta \in J\}$, then we know that:
\begin{equation*}
\forall \beta \in J, \exists \alpha'' \in I: E_\beta \in S_{\alpha''}
\end{equation*}
where the countable subset of $E$ is simply $\{E_\beta\}$ for that specific $\beta \in J$. Thus:
\begin{equation*}
E \subseteq \bigcup_{\alpha \in I}S_{\alpha} \subseteq \bigcup_{\alpha \in I}\mathcal{F}(S_\alpha) \implies \mathcal{F}(E) \subseteq \bigcup_{\alpha \in I}\mathcal{F}(S_\alpha)
\end{equation*}
where we once again applied the subset property and the fact that the $\sigma$-algebra generated by a $\sigma$-algebra is itself. Then if we take some $C \in \mathcal{F}(E)$, we have that:
\begin{equation*}
\exists \alpha \in I: C \in \mathcal{F}(S_\alpha)
\end{equation*}
So any set in the $\sigma$-algebra generated by a possibly large (uncountable), complex set can be in the simpler $\sigma$-algebra generated by a countable subset of the set.
\end{proof}

\textbf{Durrett 1.1.1:}

\textbf{Durrett 1.1.2:} Show that $\mathcal{F}(S_d) = \mathcal{B}^d$ where $\mathcal{B}$ is the Borel set of $\mathbb{R}$. Where $S_d$ is defined as:

\textbf{Durrett 1.1.3:}

\textbf{Durrett 1.1.5:}

\subsection{Measure}
The third and final element of the triple is $P$, which we call the probability measure. A measure is defined as a following:
\begin{defn}
$(X, \mathcal{F})$ is a measurable space on which we can put a measure $\mu: \mathcal{F} \xrightarrow{} [0, \infty]$ with the properties:
\begin{enumerate}
\item $\forall S \in \mathcal{F}: \mu(S) \geqslant 0 \land \mu(\varnothing) = 0$

Note that we could also equivalently replace the second part with $\mu(\varnothing) \in [0, \infty)$.
\item Given some $\{S_n\}_{n \in \mathbb{N}} \in \mathcal{F}^n$, if:
\begin{equation*}
\forall i \neq j: S_i \cap S_j = \varnothing \implies \mu\left(\bigcup_{n \in \mathbb{N}}S_n\right) = \sum_{n \in \mathbb{N}}\mu(S_n)
\end{equation*}
\end{enumerate}
This second property will be quite important, and we shall call it countable additivity. Countable additivity gives a reason as to why we wanted property $3$ in Definition 1.1, where:
\begin{equation*}
\forall \{S_n\}_{n \in \mathbb{N}} \in \mathcal{F}^\mathbb{N}: \bigcup_{n \in \mathbb{N}}S_n \in \mathcal{F}
\end{equation*}
If this were not true, then we could have had some sequence of $S_n$ in the $\sigma$-algebra where their countably infinite union is not measurable, so property $2$ in this definition would not always hold.
\end{defn}


\begin{lem}
This definition of a measure reveals additional properties:
\begin{enumerate}
\item Monotonicity:
\begin{equation*}
\forall A, B \in \mathcal{F}: A \subseteq B \implies \mu(A) \leqslant \mu(B)
\end{equation*}
\item Subadditivity:
\begin{equation*}
\forall \{S_n\}_{n \in \mathbb{N}} \in \mathcal{F}^\mathbb{N}: \mu\left(\bigcup_{n \in \mathbb{N}}S_n\right) \leqslant \sum_{n \in \mathbb{N}}\mu(S_n)
\end{equation*}
\item Continuity from below:
\begin{equation*}
\forall n \in \mathbb{N}: S_n \subseteq S_{n+1} \implies \mu\left(\bigcup_{n \in \mathbb{N}}S_n\right) = \lim_{n \xrightarrow{} \infty}\mu(S_n)
\end{equation*}
\item Continuity from above:
\begin{equation*}
\forall n \in \mathbb{N}: S_{n+1} \subseteq S_n \land \mu(S_0) < \infty \implies \mu\left(\bigcap_{n \in \mathbb{N}}S_n\right) = \lim_{n \xrightarrow{} \infty}\mu(S_n)
\end{equation*}
\end{enumerate}
\end{lem}
\begin{proof}
Take the sets $A, B \setminus A$, then:
\begin{equation*}
\mu(B) = \mu(A \cup (B \setminus A)) = \mu(A) + \mu(B \setminus A) \geqslant \mu(A)
\end{equation*}
thus proving the first property of non-negativity. Here we used countable additivity with the sequence $A, B \setminus A, \varnothing, \varnothing,...$ to bring about finite additivity. This means the proof also requires
that the measure $\mu(\varnothing) = 0$.

For the second property, we can rewrite the union as the following:
\begin{equation*}
\bigcup_{n \in \mathbb{N}}S_n = \bigcup_{n \in \mathbb{N}}\left(S_n \setminus \bigcup_{i=0}^{n-1}S_i\right)
\end{equation*}
The reverse inclusion is trivial. For the forward inclusion, if $x$ is in the union of $S_n$, then:
\begin{equation*}
\exists n \in \mathbb{N}: x \in S_n
\end{equation*}
but it might be that $x$ belongs to the intersection of multiple such sets, so by well ordering property and nonemptiness above, we can define:
\begin{equation*}
m \coloneq \min\{n \in \mathbb{N}: x \in S_n\}
\end{equation*}
So $x \in S_m$ which also tells us that for all $n \leqslant m$ that $x \notin S_n$, so:
\begin{equation*}
\exists m \in \mathbb{N}: x \in S_m \setminus \bigcup_{i=0}^{m-1}S_i
\end{equation*}
We also have the added benefit that each set difference is disjoint from any other. So:
\begin{equation*}
\mu\left(\bigcup_{n \in \mathbb{N}}S_n\right) = \mu\left(\bigcup_{n \in \mathbb{N}} \left(S_n \setminus \bigcup_{i=0}^{n-1}S_i\right)\right) = \sum_{n \in \mathbb{N}}\mu\left(S_n \setminus \bigcup_{i=0}^{n-1}S_i\right) \leqslant \sum_{n \in \mathbb{N}}\mu(S_n)
\end{equation*}
Now observe we can do the following rewrite:
\begin{equation*}
\sum_{n \in \mathbb{N}}\mu\left(S_n \setminus \bigcup_{i=0}^{n-1}S_i\right) = \lim_{n \xrightarrow{} \infty}\sum_{i=0}^{n}\mu\left(S_i \setminus \bigcup_{j=0}^{i-1}S_j\right)
\end{equation*}
Countable additivity gives us finite additivity (simply consider a sequence $A_0, A_1, A_2, ..., A_n, \varnothing, \varnothing,...$), so we can rewrite the sum as:
\begin{equation*}
= \lim_{n \xrightarrow{} \infty}\mu\left(\bigcup_{i=0}^n S_i\right) = \mu\left(\bigcup_{n \in \mathbb{N}} S_i\right)
\end{equation*}

The reason why the third and fourth properties are dubbed "continuity" is because they essentially explain a limit exchange under certain sufficient conditions, where:
\begin{equation*}
\mu\left(\bigcup_{n \in \mathbb{N}}S_i\right) = \mu\left(\lim_{n \xrightarrow{} \infty}S_n\right) = \lim_{n \xrightarrow{} \infty}\mu(S_n)
\end{equation*}
The proof follows simply from our observation before. Due to the nesting where $S_0 \subseteq S_1 \subseteq ...$, we have that:
\begin{equation*}
\bigcup_{i=0}^nS_i = S_n
\end{equation*}
By the monotonicity property, we have that $\{\mu(S_n)\}_{n \in \mathbb{N}}$ is a non-decreasing sequence, so we can rewrite the limit as:
\begin{equation*}
\lim_{n \xrightarrow{} \infty} \mu(S_n) = \sup_{n \in \mathbb{N}}\mu(S_n) 
\end{equation*}


For the fourth property, we use continuity from below. Consider a sequence $S_0 \supseteq S_1 \supseteq ...$. Then we have that $S_0 \setminus S_1 \subseteq S_0 \setminus S_2 \subseteq ...$, which tells us that:
\begin{equation*}
\mu\left(\lim_{n \xrightarrow{} \infty}S_0 \setminus S_n\right) = \lim_{n \xrightarrow{} \infty}\mu(S_0 \setminus S_n) 
\end{equation*}
Now observe:
\begin{equation*}
\mu(S_0) = \mu(S_0 \setminus S_n \cup S_n) = \mu(S_0 \setminus S_n) + \mu(S_n) \implies \mu(S_0) - \lim_{n \xrightarrow{} \infty}\mu(S_n) = \lim_{n \xrightarrow{} \infty}\mu(S_0 \setminus S_n)
\end{equation*}
\begin{equation*}
\implies \mu(S_0) - \lim_{n \xrightarrow{} \infty}\mu(S_n) = \mu(S_0 \setminus \cap S_n) = \mu(S_0) - \mu\left(\bigcap_{n \in \mathbb{N}}S_n\right)
\end{equation*}
and via cancellation we have the desired property. Notice that the argument requires $\mu(S_0) < \infty$. If it were not less than infinity, then it may be possible that even the subsets of $S_0$ have infinite measure, which would mean:
\begin{equation*}
\lim_{n \xrightarrow{} \infty}\mu(S_n) = \infty
\end{equation*}
which would give an indeterminate subtraction on the left hand side.
\end{proof}

Furthermore, we can provide an equivalent definition for measure which does not require countable additivity:
\begin{lem}
The function $\mu$ is a measure if and only if it is non-negative, finitely additive, and continuous from below.
\end{lem}
\begin{proof}
$\implies:$ This direction was proved in Lemma 1.3.


$\impliedby$: Consider a sequence of pairwise disjoint sets $A_1, A_2, ...$. Then we have the sequence:
\begin{equation*}
\left\{\bigcup_{i=0}^n A_n\right\}_{n \in \mathbb{N}}
\end{equation*}
which is non-decreasing (satisfying our conditiion for continuity from below). Then:
\begin{equation*}
\mu\left(\bigcup_{n \in \mathbb{N}}A_n\right) = \mu\left(\bigcup_{n \in \mathbb{N}}\left(\bigcup_{i=0}^n A_n\right)\right) = \lim_{n \xrightarrow{} \infty}\mu\left(\bigcup_{i=0}^n A_n\right)
\end{equation*}
By finite additivity, because the sequence is pairwise disjoint, we have that:
\begin{equation*}
= \lim_{n \xrightarrow{} \infty}\sum_{i=0}^{n}\mu(A_i) = \sum_{n \in \mathbb{N}}\mu(A_n)
\end{equation*}
Now we need to prove $\mu(\varnothing) = 0$. By finite additivity and $\mu(\varnothing) < \infty$:
\begin{equation*}
\mu(\varnothing \cup \varnothing) = \mu(\varnothing) = \mu(\varnothing) + \mu(\varnothing) \implies \mu(\varnothing) = 0
\end{equation*}
So we must assume non-negativity (if we try to prove it via monotonicity, it becomes circular). Then we have that:
\begin{equation*}
(\mu(\varnothing) \in [0, \infty)) + \text{finitely additive + continuous from below} \iff (\mu(\varnothing) = 0)  + \text{countably additive}
\end{equation*}
\begin{equation*}
\iff (\mu(\varnothing) \in [0, \infty)) + \text{countably additive}
\end{equation*}

\end{proof}
We can equivalently define a measure as:
\begin{defn}[Alternative Definition of Measure]
A measure $\mu: \mathcal{F} \xrightarrow{} [0, \infty]$ has the properties:
\begin{enumerate}
\item For all finite collections of pairwise disjoint sets $S_0, S_1, ..., S_n$:
\begin{equation*}
\mu\left(\bigcup_{i=0}^n S_i\right) = \sum_{i=0}^{n}\mu(S_i)
\end{equation*}
\item For all sequences $\{S_n\}_{n \in \mathbb{N}} \in \mathcal{F}^\mathbb{N}$:
\begin{equation*}
\forall n \in \mathbb{N}: S_n \subseteq S_{n+1} \implies \mu\left(\bigcup_{n \in \mathbb{N}}S_n\right) = \lim_{n \xrightarrow{} \infty}\mu(S_n)
\end{equation*}
\end{enumerate}
\end{defn}
We can now make this notion more specific to probability and say:
\begin{defn}
A probability measure is a function $P: \mathcal{F} \xrightarrow{} [0, \infty]$ where:
\begin{enumerate}
\item $P(\Omega) = 1$
\item Given some sequence $\{S_n\}_{n \in \mathbb{N}} \in \mathcal{F}^\mathbb{N}$:
\begin{equation*}
\forall i \neq j: S_i \cap S_j = \varnothing \implies P\left(\bigcup_{n \in \mathbb{N}}S_n\right) = \sum_{n \in \mathbb{N}}P(S_n)
\end{equation*}
\end{enumerate}
\end{defn}

\begin{lem}
$P$ is a measure on the measurable space $(\Omega, \mathcal{F})$ and:
\begin{equation*}
\forall S \in \mathcal{F}: P(\Omega \setminus S) = 1 - P(S)
\end{equation*}
\end{lem}
\begin{proof}
Non-negativity and countable additivity follows by definition. We only need to prove that $\mu(\varnothing) = 0$. Simply take the sequence:
\begin{equation*}
\Omega, \varnothing, \varnothing,...
\end{equation*}
Then by countable additivity:
\begin{equation*}
P(\Omega) = P(\Omega) + \sum_{n \in \mathbb{N}}P(\varnothing)
\end{equation*}
Since $P(\Omega) = 1$, we can apply cancellation, which forces $P(\varnothing) = 0$. 

This is a sufficient condition for measure spaces in general. Given a non-negative function $\mu$ with countable additivity:
\begin{equation*}
\exists S \in \mathcal{F}: (\mu(S) < \infty) + \text{countable additivity} \implies \mu(\varnothing) = 0
\end{equation*}
proving that $\mu$ is a measure on that measurable space.

We prove the second part now that we know that $P(\varnothing) = 0$. Because of this fact, we can say that countable additivity implies finite additivity, thus:
\begin{equation*}
P(\Omega) = P(\Omega \setminus S \cup S) = P(\Omega \setminus S) + P(S)
\end{equation*}
\begin{equation*}
\implies P(\Omega \setminus S) = 1 - P(S)
\end{equation*}
To do this proof, we require that $S \subseteq \Omega$ for the decomposition to hold.
\end{proof}

\begin{defn}
If 
\end{defn}



To motivate the use of $\sigma$-algebras, we look at a simple probability based example. Consider a weighted coin where the conductor of the experiment is not aware of the weightage. Then:
\begin{equation*}
\Omega \coloneq \{H, T\}
\end{equation*}
However, our $\sigma$-algebra would be:
\begin{equation*}
\mathcal{F} \coloneq \{\Omega, \varnothing\}
\end{equation*}
because $P(\{H\})$ or $P(\{T\})$ would not be known. We can make a more complicated example.


we fix our attention to subsets of $\mathbb{R}$. We defined the notion of outer measure of subsets of $\mathbb{R}$ as the following:
\begin{equation*}
\forall A \subseteq \mathbb{R}: |A| \coloneq \inf\left\{\sum_{n \in \mathbb{N}}|I_n| : A \subseteq \bigcup_{n \in \mathbb{N}}I_n\right\}
\end{equation*}
where $I_n$ is an open interval (or closed works as well) whose measure we defined as the difference between its endpoints $b-a$.
\pagebreak

\section{Random Variables and Distributions}

\end{document}